
\documentclass[12pt,a4paper]{article}

\usepackage[brazilian]{babel}

\usepackage{amsmath}
\usepackage{amssymb}
\usepackage{enumitem}
\usepackage{etoolbox}
\usepackage{tocloft}
\usepackage{tikz}

\setlength{\cftsecnumwidth}{3.2cm}
\setlength{\cftsubsecnumwidth}{1.5cm}
\setlength{\cftsubsecnumwidth}{1.5cm}

\renewcommand{\thesection}{Capítulo \arabic{section} -}
\renewcommand{\thesubsection}{\arabic{section}.\arabic{subsection}}
\pretocmd{\subsection}{\clearpage}{}{}

\title{
    \Huge \textbf{Um Curso de Cálculo Volume 1} \\
    \vspace{0.3cm}
    \Large 6\textordfeminine{} Edição (Guidorizzi) \\
    \vspace{0.5cm}
    \large \textit{Resoluções Completas e Passo a Passo}
}
\author{Rogério Campos}

\begin{document}

\maketitle
\newpage

\tableofcontents
\newpage

\setcounter{section}{6}
\setcounter{subsection}{1}

\section{Derivadas}

\subsection{Introdução}
Esta subseção não possui exercícios

\subsection{Derivada de uma função}
\vspace{0.5cm}
\noindent\textbf{Exercício 1.} Seja $f(x) = x^2 + 1$. Calcule:

\begin{enumerate}
    \renewcommand{\theenumi}{\alph{enumi}}
    \renewcommand{\labelenumi}{\theenumi)}
    \item $f'(x)$
    \item $f'(0)$
    \item $f'(1)$
\end{enumerate}

\vspace{0.5cm}
\noindent\textbf{Exercício 2.} Seja $f(x) = 2x$. Pensando geometricamente, qual o valor que você espera para $f'(p)$? Calcule $f'(p)$.

\vspace{0.5cm}
\noindent\textbf{Exercício 3.} Seja $f(x) = 3x + 2$. Calcule:

\begin{enumerate}
    \renewcommand{\theenumi}{\alph{enumi}}
    \renewcommand{\labelenumi}{\theenumi)}
    \item $f'(2)$
    \item $f'(0)$
    \item $f'(x)$
\end{enumerate}

\vspace{0.5cm}
\noindent\textbf{Exercício 4.} Calcule $f'(p)$, pela definição, sendo dados:

\begin{enumerate}
    \renewcommand{\theenumi}{\alph{enumi}}
    \renewcommand{\labelenumi}{\theenumi)}
    \item $f(x) = x^2 + x$ e $p = 1$
    \item $f(x) = \sqrt{x}$ e $p = 4$
    \item $f(x) = 5x - 3$ e $p = -3$
    \item $f(x) = \frac{1}{x}$ e $p = 1$
    \item $f(x) = \sqrt{x}$ e $p = 3$
    \item $f(x) = \frac{1}{\sqrt{x^2}}$ e $p = 2$
    \item $f(x) = 2x^3 - x^2$ e $p = 1$
    \item $f(x) = \sqrt[3]{x}$ e $p = 2$
\end{enumerate}


\vspace{0.5cm}
\noindent\textbf{Exercício 5.} Determine a equação da reta tangente em $(p, f(p))$ sendo dados:

\begin{enumerate}
    \renewcommand{\theenumi}{\alph{enumi}}
    \renewcommand{\labelenumi}{\theenumi)}
    \item $f(x) = x^2$ e $p = 2$
    \item $f(x) = \frac{1}{x}$ e $p = 2$
    \item $f(x) = \sqrt{x}$ e $p = 9$
    \item $f(x) = x^2 - x$ e $p = 1$
\end{enumerate}


\vspace{0.5cm}
\noindent\textbf{Exercício 6.} Calcule $f'(x)$, pela definição:

\begin{enumerate}
    \renewcommand{\theenumi}{\alph{enumi}}
    \renewcommand{\labelenumi}{\theenumi)}
    \item $f(x) = x^2 + x$
    \item $f(x) = 3x - 1$
    \item $f(x) = x^3$
    \item $f(x) = \frac{1}{x}$
    \item $f(x) = 5x$
    \item $f(x) = 10$
    \item $f(x) = \frac{x}{x + 1}$
    \item $f(x) = \frac{1}{x^2}$
\end{enumerate}

\vspace{0.5cm}
\noindent\textbf{Exercício 7.} Dê exemplo (por meio de um gráfico) de uma função $f$, definida e derivável em $\mathbb{R}$, tal que $f'(1) = 0$.

\vspace{0.5cm}
\noindent\textbf{Exercício 8.} Dê exemplo (por meio de um gráfico) de uma função $f$, definida e derivável em $\mathbb{R}$, tal que $f'(x) > 0$ para todo $x$.

\vspace{0.5cm}
\noindent\textbf{Exercício 9.} Dê exemplo (por meio de um gráfico) de uma função $f$, definida e derivável em $\mathbb{R}$, tal que $f'(0) < f'(1)$.

\vspace{0.5cm}
\noindent\textbf{Exercício 10.} Dê exemplo (por meio de um gráfico) de uma função $f$, definida e contínua em $\mathbb{R}$, tal que $f'(1)$ não exista.

\vspace{0.5cm}
\noindent\textbf{Exercício 11.} Dê exemplo (por meio de um gráfico) de uma função $f$, definida e derivável em $\mathbb{R}$, tal que $f'(x) > 0$ para $x < 1$ e $f'(x) < 0$ para $x > 1$.

\vspace{0.5cm}
\noindent\textbf{Exercício 12.} Dê exemplo (por meio de um gráfico) de uma função $f$, definida e derivável em $\mathbb{R}$, tal que $f'(x) > 0$ para $x < 0$, $f'(x) < 0$ para $0 < x < 2$ e $f'(x) > 0$ para $x > 2$.

\vspace{0.5cm}
\noindent\textbf{Exercício 13.} Dê exemplo (por meio de um gráfico) de uma função $f$, definida e derivável em $\mathbb{R}$, tal que $f'(0) = 0$ e $f'(1) = 0$.

\vspace{0.5cm}
\noindent\textbf{Exercício 14.} Mostre que a função
\[
g(x) = \begin{cases} 
2x + 1 & \text{se } x < 1 \\ 
-x + 4 & \text{se } x \geq 1 
\end{cases}
\]
não é derivável em $p = 1$. Esboce o gráfico de $g$.

\vspace{0.5cm}
\noindent\textbf{Exercício 15.} Seja $g(x) = \begin{cases} x^2 + 2 & \text{se } x < 1 \\ 2x + 1 & \text{se } x \geq 1 \end{cases}$
\begin{enumerate}
    \renewcommand{\theenumi}{\alph{enumi}}
    \renewcommand{\labelenumi}{\theenumi)}
    \item Mostre que $g$ é derivável em $p = 1$ e calcule $g'(1)$.
    \item Esboce o gráfico de $g$.
\end{enumerate}

\vspace{0.5cm}
\noindent\textbf{Exercício 16.} Seja $f(x) = \begin{cases} 2 & \text{se } x \geq 0 \\ x^2 + 2 & \text{se } x < 0 \end{cases}$
\begin{enumerate}
    \renewcommand{\theenumi}{\alph{enumi}}
    \renewcommand{\labelenumi}{\theenumi)}
    \item Esboce o gráfico de $f$.
    \item $f$ é derivável em $p = 0$? Em caso afirmativo, calcule $f'(0)$.
\end{enumerate}

\vspace{0.5cm}
\noindent\textbf{Exercício 17.} Seja $g(x) = \begin{cases} x + 1 & \text{se } x < 1 \\ -x + 3 & \text{se } x \geq 1 \end{cases}$
\begin{enumerate}
    \renewcommand{\theenumi}{\alph{enumi}}
    \renewcommand{\labelenumi}{\theenumi)}
    \item Esboce o gráfico de $g$.
    \item $g$ é derivável em $p = 1$? Por quê?
\end{enumerate}

\vspace{0.5cm}
\noindent\textbf{Exercício 18.}
Construa uma função $f\colon \mathbb{R} \to \mathbb{R}$ que seja contínua em $\mathbb{R}$ e que seja derivável em todos os pontos, exceto em $-1$, $0$ e $1$.

\vspace{0.5cm}
\noindent\textbf{Exercício 19.}
Construa uma função $f\colon \mathbb{R} \to \mathbb{R}$ que seja contínua em $\mathbb{R}$ e derivável em todos os pontos, exceto nos números inteiros.

\subsection{Derivadas de $x^n$ e $\sqrt[n]{x}$}
\vspace{0.5cm}
\noindent\textbf{Exercício 1.} Seja $f(x) = x^5$. Calcule
\begin{enumerate}
    \renewcommand{\theenumi}{\alph{enumi}}
    \renewcommand{\labelenumi}{\theenumi)}
    \item $f'(x)$
    \item $f'(0)$
    \item $f'(2)$
\end{enumerate}

\vspace{0.5cm}
\noindent\textbf{Exercício 2.} Calcule $g'(x)$ sendo $g$ dada por
\begin{enumerate}
    \renewcommand{\theenumi}{\alph{enumi}}
    \renewcommand{\labelenumi}{\theenumi)}
    \item $g(x) = x^6$
    \item $g(x) = x^{100}$
    \item $g(x) = \dfrac{1}{x}$
    \item $g(x) = x^2$
    \item $g(x) = \dfrac{1}{x^3}$
    \item $g(x) = \dfrac{1}{x^7}$
    \item $g(x) = x$
    \item $g(x) = x^{-3}$
\end{enumerate}

\vspace{0.5cm}
\noindent\textbf{Exercício 3.}
Determine a equação da reta tangente ao gráfico de $f(x) = \dfrac{1}{x}$ no ponto de abscissa 2. Esboce os gráficos de $f$ e da reta tangente.

\vspace{0.5cm}
\noindent\textbf{Exercício 4.}
Determine a equação da reta tangente ao gráfico de $f(x) = \dfrac{1}{x^2}$ no ponto de abscissa 1. Esboce os gráficos de $f$ e da reta tangente.

\vspace{0.5cm}
\noindent\textbf{Exercício 5.} Seja $f(x) = \sqrt[5]{x}$. Calcule.
\begin{enumerate}
    \renewcommand{\theenumi}{\alph{enumi}}
    \renewcommand{\labelenumi}{\theenumi)}
    \item $f'(x)$
    \item $f'(1)$
    \item $f'(-32)$
\end{enumerate}

\vspace{0.5cm}
\noindent\textbf{Exercício 6.} Calcule $g'(x)$, sendo $g$ dada por
\begin{enumerate}
    \renewcommand{\theenumi}{\alph{enumi}}
    \renewcommand{\labelenumi}{\theenumi)}
    \item $g(x) = \sqrt[4]{x}$
    \item $g(x) = \sqrt[6]{x}$
    \item $g(x) = \sqrt[8]{x}$
    \item $g(x) = \sqrt[9]{x}$
\end{enumerate}

\vspace{0.5cm}
\noindent\textbf{Exercício 7.}
Determine a equação da reta tangente ao gráfico de $f(x) = \sqrt[3]{x}$ no ponto de abscissa 1. Esboce os gráficos de $f$ e da reta tangente.

\vspace{0.5cm}
\noindent\textbf{Exercício 8.}
Seja $r$ a reta tangente ao gráfico de $f(x) = \dfrac{1}{x}$ no ponto de abscissa $p$. Verifique que $r$ intercepta o eixo $x$ no ponto de abscissa $2p$.

\vspace{0.5cm}
\noindent\textbf{Exercício 9.}
Determine a reta que é tangente ao gráfico de $f(x) = x^2$ e paralela à reta $y = 4x + 2$.

\subsection{Derivadas de $e^x$ e $\ln(x)$}
\vspace{0.5cm}
\noindent\textbf{Exercício 1.}
Determine a equação da reta tangente ao gráfico de $f(x) = e^x$ no ponto de abscissa 0.

\vspace{0.5cm}
\noindent\textbf{Exercício 2.}
Determine a equação da reta tangente ao gráfico de $f(x) = \ln x$ no ponto de abscissa 1. Esboce os gráficos de $f$ e da reta tangente.

\vspace{0.5cm}
\noindent\textbf{Exercício 3.}
Seja $f(x) = a^x$, em que $a > 0$ e $a \neq 1$ é um real dado. Mostre que $f'(x) = a^x \ln a$.

\vspace{0.5cm}
\noindent\textbf{Exercício 4.} Calcule $f'(x)$.
\begin{enumerate}
    \renewcommand{\theenumi}{\alph{enumi}}
    \renewcommand{\labelenumi}{\theenumi)}
    \item $f(x) = 2^x$
    \item $f(x) = 5^x$
    \item $f(x) = \pi^x$
    \item $f(x) = e^x$
\end{enumerate}

\vspace{0.5cm}
\noindent\textbf{Exercício 5.}
Seja $g(x) = \log_a x$, em que $a > 0$ e $a \neq 1$ é constante. Mostre que $g'(x) = \dfrac{1}{x \ln a}$.

\vspace{0.5cm}
\noindent\textbf{Exercício 6.} Calcule $g'(x)$
\begin{enumerate}
    \renewcommand{\theenumi}{\alph{enumi}}
    \renewcommand{\labelenumi}{\theenumi)}
    \item $g(x) = \log_3 x$
    \item $g(x) = \log_5 x$
    \item $g(x) = \log_\pi x$
    \item $g(x) = \ln x$
\end{enumerate}

\subsection{Derivadas de Funções Trigonométricas}
\vspace{0.5cm}
\noindent\textbf{Exercício 1.} Seja $f(x) = \operatorname{sen} x$. Calcule.
\begin{enumerate}
    \renewcommand{\theenumi}{\alph{enumi}}
    \renewcommand{\labelenumi}{\theenumi)}
    \item $f'(x)$
    \item $f'\!\left(\dfrac{\pi}{4}\right)$
\end{enumerate}

\vspace{0.5cm}
\noindent\textbf{Exercício 2.}
Determine a equação da reta tangente ao gráfico de $f(x) = \operatorname{sen} x$ no ponto de abscissa 0.

\vspace{0.5cm}
\noindent\textbf{Exercício 3.} Seja $f(x) = \cos x$. Calcule.
\begin{enumerate}
    \renewcommand{\theenumi}{\alph{enumi}}
    \renewcommand{\labelenumi}{\theenumi)}
    \item $f'(x)$
    \item $f'(0)$
    \item $f'\!\left(\dfrac{\pi}{3}\right)$
    \item $f'\!\left(-\dfrac{\pi}{4}\right)$
\end{enumerate}

\vspace{0.5cm}
\noindent\textbf{Exercício 4.} Calcule $f'(x)$ sendo
\begin{enumerate}
    \renewcommand{\theenumi}{\alph{enumi}}
    \renewcommand{\labelenumi}{\theenumi)}
    \item $f(x) = \operatorname{tg} x$
    \item $f(x) = \sec x$
\end{enumerate}

\vspace{0.5cm}
\noindent\textbf{Exercício 5.}
Determine a equação da reta tangente ao gráfico de $f(x) = \operatorname{tg} x$ no ponto de abscissa 0.

\vspace{0.5cm}
\noindent\textbf{Exercício 6.} Seja $f(x) = \operatorname{cotg} x$. Calcule.
\begin{enumerate}
    \renewcommand{\theenumi}{\alph{enumi}}
    \renewcommand{\labelenumi}{\theenumi)}
    \item $f'(x)$
    \item $f'\!\left(\dfrac{\pi}{4}\right)$
\end{enumerate}

\vspace{0.5cm}
\noindent\textbf{Exercício 7.} Seja $g(x) = \operatorname{cosec} x$. Calcule.
\begin{enumerate}
    \renewcommand{\theenumi}{\alph{enumi}}
    \renewcommand{\labelenumi}{\theenumi)}
    \item $g'(x)$
    \item $g'\!\left(\dfrac{\pi}{4}\right)$
\end{enumerate}

\subsection{Derivabilidade e Continuidade}
\vspace{0.5cm}
\noindent\textbf{Exercício 1.} Seja $f(x) = \begin{cases} x + 1 & \text{se } x < 2 \\ 1 & \text{se } x \ge 2 \end{cases}$
\begin{enumerate}
    \renewcommand{\theenumi}{\alph{enumi}}
    \renewcommand{\labelenumi}{\theenumi)}
    \item $f$ é contínua em 2? Por quê?
    \item $f$ é derivável em 2? Por quê?
\end{enumerate}

\vspace{0.5cm}
\noindent\textbf{Exercício 2.} Seja $f(x) = \begin{cases} x^2 & \text{se } x \le 0 \\ -x^2 & \text{se } x > 0 \end{cases}$
\begin{enumerate}
    \renewcommand{\theenumi}{\alph{enumi}}
    \renewcommand{\labelenumi}{\theenumi)}
    \item $f$ é derivável em 0? Justifique.
    \item $f$ é contínua em 0? Justifique.
\end{enumerate}

\vspace{0.5cm}
\noindent\textbf{Exercício 3.} Seja $f(x) = \begin{cases} -x + 3 & \text{se } x < 3 \\ x - 3 & \text{se } x \ge 3 \end{cases}$
\begin{enumerate}
    \renewcommand{\theenumi}{\alph{enumi}}
    \renewcommand{\labelenumi}{\theenumi)}
    \item $f$ é derivável em 3? Justifique.
    \item $f$ é contínua em 3? Justifique.
\end{enumerate}

\subsection{Regras de Derivação}
\vspace{0.5cm}
\noindent\textbf{Exercício 1.} Calcule $f'(x)$.
\begin{enumerate}
    \renewcommand{\theenumi}{\alph{enumi}}
    \renewcommand{\labelenumi}{\theenumi)}
    \item $f(x) = 3x^2 + 5$
    \item $f(x) = x^3 + x^2 + 1$
    \item $f(x) = 3x^3 - 2x^2 + 4$
    \item $f(x) = 3x + \sqrt{x}$
    \item $f(x) = 5 + 3x^{-2}$
    \item $f(x) = 2\sqrt[3]{x}$
    \item $f(x) = 3x + \dfrac{1}{x}$
    \item $f(x) = \dfrac{4}{x} + \dfrac{5}{x^2}$
    \item $f(x) = \dfrac{2}{3}x^3 + \dfrac{1}{4}x^2$
    \item $f(x) = \sqrt[3]{x} + \sqrt{x}$
    \item $f(x) = 2x + \dfrac{1}{x} + \dfrac{1}{x^2}$
    \item $f(x) = 6x^3 + \sqrt[3]{x}$
    \item $f(x) = 5x^4 + bx^3 + cx^2 + k$, em que $b$, $c$ e $k$ são constantes.
\end{enumerate}

\vspace{0.5cm}
\noindent\textbf{Exercício 2.}
Seja $g(x) = x^3 + \dfrac{1}{x}$. Determine a equação da reta tangente ao gráfico de $g$ no ponto $(1, g(1))$.

\vspace{0.5cm}
\noindent\textbf{Exercício 3.} Seja $f(x) = x^2 + \dfrac{1}{x}$.
\begin{enumerate}
    \renewcommand{\theenumi}{\alph{enumi}}
    \renewcommand{\labelenumi}{\theenumi)}
    \item Determine o ponto do gráfico de $f$ em que a reta tangente, neste ponto, seja paralela ao eixo $x$.
    \item Esboce o gráfico de $f$.
\end{enumerate}

\vspace{0.5cm}
\noindent\textbf{Exercício 4.} Seja $f(x) = x^3 + 3x^2 + 1$.
\begin{enumerate}
    \renewcommand{\theenumi}{\alph{enumi}}
    \renewcommand{\labelenumi}{\theenumi)}
    \item Estude o sinal de $f'(x)$.
    \item Calcule $\lim\limits_{x \to +\infty} f(x)$ e $\lim\limits_{x \to -\infty} f(x)$.
    \item Utilizando as informações acima, faça um esboço do gráfico de $f$.
\end{enumerate}

\vspace{0.5cm}
\noindent\textbf{Exercício 5.}
Mesmo exercício que o anterior, considerando a função $f(x) = x^3 + x^2 - 5x$.

\vspace{0.5cm}
\noindent\textbf{Exercício 6.} Seja $f(x) = x^3 + 3x$.
\begin{enumerate}
    \renewcommand{\theenumi}{\alph{enumi}}
    \renewcommand{\labelenumi}{\theenumi)}
    \item Determine a equação da reta tangente ao gráfico de $f$ no ponto de abscissa 0.
    \item Estude o sinal de $f'(x)$.
    \item Esboce o gráfico de $f$.
\end{enumerate}

\vspace{0.5cm}
\noindent\textbf{Exercício 7.} Calcule $F'(x)$ em que $F(x)$ é igual a
\begin{enumerate}
    \renewcommand{\theenumi}{\alph{enumi}}
    \renewcommand{\labelenumi}{\theenumi)}
    \item $\dfrac{x}{x^2 + 1}$
    \item $\dfrac{x^2 - 1}{x + 1}$
    \item $\dfrac{3x^2 + 3}{5x - 3}$
    \item $\dfrac{\sqrt{x}}{x + 1}$
    \item $5x + \dfrac{x}{x - 1}$
    \item $\sqrt{x} + \dfrac{3}{x^3 + 2}$
    \item $\dfrac{\sqrt[3]{x} + x}{\sqrt{x}}$
    \item $\dfrac{x + \sqrt[4]{x}}{x^2 + 3}$
\end{enumerate}

\vspace{0.5cm}
\noindent\textbf{Exercício 8.} Seja $g(x) = \dfrac{x}{x^2 + 1}$.
\begin{enumerate}
    \renewcommand{\theenumi}{\alph{enumi}}
    \renewcommand{\labelenumi}{\theenumi)}
    \item Determine os pontos do gráfico de $g$ em que as retas tangentes, nestes pontos, sejam paralelas ao eixo $x$.
    \item Estude o sinal de $g'(x)$.
    \item Calcule $\lim\limits_{x \to +\infty} g(x)$ e $\lim\limits_{x \to -\infty} g(x)$.
    \item Utilizando as informações acima, faça um esboço do gráfico de $g$.
\end{enumerate}

\vspace{0.5cm}
\noindent\textbf{Exercício 9.} Calcule $f'(x)$ em que $f(x)$ é igual a
\begin{enumerate}
    \renewcommand{\theenumi}{\alph{enumi}}
    \renewcommand{\labelenumi}{\theenumi)}
    \item $3x^2 + 5\cos x$
    \item $\dfrac{\cos x}{x^2 + 1}$
    \item $x\operatorname{sen} x$
    \item $x^2 \operatorname{tg} x$
    \item $\dfrac{x + 1}{\operatorname{tg} x}$
    \item $\dfrac{3}{\operatorname{sen} x + \cos x}$
    \item $\dfrac{\sec x}{3x + 2}$
    \item $\cos x + (x^2 + 1)\operatorname{sen} x$
    \item $\sqrt{x}\sec x$
    \item $3\cos x + 5\sec x$
    \item $x\operatorname{cotg} x$
    \item $4\sec x + \operatorname{cotg} x$
    \item $x^2 + 3x\operatorname{tg} x$
    \item $\dfrac{x^2 + 1}{\sec x}$
    \item $\dfrac{x + 1}{x\operatorname{sen} x}$
    \item $\dfrac{x}{\operatorname{cosec} x}$
    \item $(x^3 + \sqrt{x})\operatorname{cosec} x$
    \item $\dfrac{x + \operatorname{sen} x}{x - \cos x}$
\end{enumerate}

\vspace{0.5cm}
\noindent\textbf{Exercício 10.} Seja $f(x) = x^2 \operatorname{sen} x + \cos x$. Calcule:
\begin{enumerate}
    \renewcommand{\theenumi}{\alph{enumi}}
    \renewcommand{\labelenumi}{\theenumi)}
    \item $f'(x)$
    \item $f'(0)$
    \item $f'(3a)$
    \item $f'(x^2)$
\end{enumerate}

\vspace{0.5cm}
\noindent\textbf{Exercício 11.} Seja $f(x) = \operatorname{sen} x + \cos x$, $0 \le x \le 2\pi$.
\begin{enumerate}
    \renewcommand{\theenumi}{\alph{enumi}}
    \renewcommand{\labelenumi}{\theenumi)}
    \item Estude o sinal de $f'(x)$.
    \item Faça um esboço do gráfico de $f$.
\end{enumerate}

\vspace{0.5cm}
\noindent\textbf{Exercício 12.} Calcule $f'(x)$.
\begin{enumerate}
    \renewcommand{\theenumi}{\alph{enumi}}
    \renewcommand{\labelenumi}{\theenumi)}
    \item $f(x) = x^2 e^x$
    \item $f(x) = 3x + 5\ln x$
    \item $f(x) = e^x \cos x$
    \item $f(x) = \dfrac{1 + e^x}{1 - e^x}$
    \item $f(x) = x^2 \ln x + 2e^x$
    \item $f(x) = \dfrac{x + 1}{x\ln x}$
    \item $f(x) = 4 + 5x^2 \ln x$
    \item $f(x) = \dfrac{e^x}{x^2 + 1}$
    \item $f(x) = \dfrac{\ln x}{x}$
    \item $f(x) = \dfrac{e^x}{x + 1}$
\end{enumerate}

\vspace{0.5cm}
\noindent\textbf{Exercício 13.}
Sejam $f$, $g$ e $h$ funções deriváveis. Verifique que
\[
[f(x)g(x)h(x)]' = f'(x)g(x)h(x) + f(x)g'(x)h(x) + f(x)g(x)h'(x).
\]

\vspace{0.5cm}
\noindent\textbf{Exercício 14.} Calcule $F'(x)$ sendo $F(x)$ igual a
\begin{enumerate}
    \renewcommand{\theenumi}{\alph{enumi}}
    \renewcommand{\labelenumi}{\theenumi)}
    \item $x e^x \cos x$
    \item $x^2 (\cos x)(1 + \ln x)$
    \item $e^x \operatorname{sen} x \cos x$
    \item $(1 + \sqrt{x})e^x \operatorname{tg} x$
\end{enumerate}

\subsection{Função Derivada e Derivadas de Ordem Superior}
\vspace{0.5cm}
\noindent\textbf{Exercício 1.} Determine $f'$, $f''$ e $f'''$.
\begin{enumerate}
    \renewcommand{\theenumi}{\alph{enumi}}
    \renewcommand{\labelenumi}{\theenumi)}
    \item $f(x) = 4x^4 + 2x$
    \item $f(x) = \dfrac{1}{x}$
    \item $f(x) = 5x^2 - \dfrac{1}{x^3}$
    \item $f(x) = 3x^3 - 6x + 1$
    \item $f(x) = x|x|$
    \item $f(x) = \begin{cases} x^2 + 3x & \text{se } x \le 1 \\ 5x - 1 & \text{se } x > 1 \end{cases}$
\end{enumerate}

\vspace{0.5cm}
\noindent\textbf{Exercício 2.} Esboce os gráficos de $f$, $f'$ e $f''$.
\begin{enumerate}
    \renewcommand{\theenumi}{\alph{enumi}}
    \renewcommand{\labelenumi}{\theenumi)}
    \item $f(x) = x^2 |x|$
    \item $f(x) = \begin{cases} x^2 + 3x & \text{se } x \le 1 \\ 5x - 1 & \text{se } x > 1 \end{cases}$
\end{enumerate}

\vspace{0.5cm}
\noindent\textbf{Exercício 3.} Determine a derivada de ordem $n$.
\begin{enumerate}
    \renewcommand{\theenumi}{\alph{enumi}}
    \renewcommand{\labelenumi}{\theenumi)}
    \item $f(x) = e^x$
    \item $f(x) = \operatorname{sen} x$
    \item $f(x) = \cos x$
    \item $f(x) = \ln x$
\end{enumerate}

\subsection{Notações para a Derivada}
\vspace{0.5cm}
\noindent\textbf{Exercício 1.} Calcule a derivada.
\begin{enumerate}
    \renewcommand{\theenumi}{\alph{enumi}}
    \renewcommand{\labelenumi}{\theenumi)}
    \item $y = 5x^3 + 6x - 1$
    \item $s = \sqrt[5]{t} + \dfrac{3}{t}$
    \item $x = \dfrac{t}{t + 1}$
    \item $y = t \cos t$
    \item $y = \dfrac{u + 1}{\ln u}$
    \item $x = t^3 e^t$
    \item $s = e^t \operatorname{tg} t$
    \item $y = \dfrac{x^3 + 1}{\operatorname{sen} x}$
    \item $y = \sqrt[3]{u} \sec u$
    \item $x = \dfrac{3}{t} + \dfrac{2}{t^2}$
    \item $x = e^t \cos t$
    \item $u = 5v^2 + \dfrac{3}{v^4}$
    \item $V = \dfrac{4}{3}\pi r^3$
    \item $E = \dfrac{1}{2}v^2$
    \item $E = \dfrac{1}{2}mv^2$, $m$ constante
    \item $U = \dfrac{a}{x^{12}} - \dfrac{b}{x^6}$, $a$ e $b$ constantes
\end{enumerate}

\vspace{0.5cm}
\noindent\textbf{Exercício 2.} Seja $y = \dfrac{x^3}{x + \sqrt{x}}$. Calcule.
\begin{enumerate}
    \renewcommand{\theenumi}{\alph{enumi}}
    \renewcommand{\labelenumi}{\theenumi)}
    \item $\dfrac{dy}{dx}$
    \item $\left.\dfrac{dy}{dx}\right|_{x = 1}$
\end{enumerate}

\vspace{0.5cm}
\noindent\textbf{Exercício 3.}
Seja $y = t^2 x$, em que $x = x(t)$ é uma função derivável. Calcule $\left.\dfrac{dy}{dt}\right|_{t = 1}$ supondo $\left.\dfrac{dx}{dt}\right|_{t = 1} = 2$ e $x = 3$ para $t = 1$ (isto é, $x(1) = 3$).

\vspace{0.5cm}
\noindent\textbf{Exercício 4.}
Considere a função $y = xt^3$, na qual $x = x(t)$ é uma função derivável. Calcule $\left.\dfrac{dy}{dt}\right|_{t = 2}$ sabendo que $\left.\dfrac{dx}{dt}\right|_{t = 2} = 3$ e que $x(2) = 1$ (isto é, $x = 1$ para $t = 2$).

\vspace{0.5cm}
\noindent\textbf{Exercício 5.}
Considere a função $y = \dfrac{t}{x + t}$, na qual $t = t(x)$ é uma função derivável. Calcule $\left.\dfrac{dy}{dx}\right|_{x = 1}$ sabendo que $\left.\dfrac{dt}{dx}\right|_{x = 1} = 4$ e que $t = 2$ para $x = 1$. (Observe que $t$ está sendo olhado como função de $x$.)

\vspace{0.5cm}
\noindent\textbf{Exercício 6.}
Seja $y = \dfrac{1}{x^2}$. Verifique que $x\dfrac{dy}{dx} + 2y = 0$.

\vspace{0.5cm}
\noindent\textbf{Exercício 7.}
Seja $y = \dfrac{-2}{x^2 + k}$, $k$ constante. Verifique que $\dfrac{dy}{dx} - xy^2 = 0$.

\vspace{0.5cm}
\noindent\textbf{Exercício 8.} Calcule a derivada segunda.
\begin{enumerate}
    \renewcommand{\theenumi}{\alph{enumi}}
    \renewcommand{\labelenumi}{\theenumi)}
    \item $y = x^3 + 2x - 3$
    \item $x = t \operatorname{sen} t$
    \item $y = x^{10} + \dfrac{1}{x^3}$
    \item $y = t \ln t$
    \item $x = e^t \cos t$
    \item $y = \dfrac{e^x}{x}$
\end{enumerate}

\vspace{0.5cm}
\noindent\textbf{Exercício 9.}
Seja $y = x^2 - 3x$. Verifique que $x \dfrac{d^2 y}{dx^2} - \dfrac{dy}{dx} = 3$.

\vspace{0.5cm}
\noindent\textbf{Exercício 10.}
Seja $y = \dfrac{1}{x}$. Verifique que $x^2 \dfrac{d^3 y}{dx^3} = 6 \dfrac{dy}{dx}$.

\vspace{0.5cm}
\noindent\textbf{Exercício 11.}
Seja $x = \cos t$. Verifique que $\dfrac{d^2 y}{dt^2} + x = 0$.

\vspace{0.5cm}
\noindent\textbf{Exercício 12.}
Seja $y = e^x \cos x$. Verifique que $\dfrac{d^2 y}{dx^2} - 2\dfrac{dy}{dx} + 2y = 0$.

\vspace{0.5cm}
\noindent\textbf{Exercício 13.}
Seja $y = te^t$. Verifique que $\dfrac{d^2 y}{dt^2} - 2\dfrac{dy}{dx} + y = 0$.

\vspace{0.5cm}
\noindent\textbf{Exercício 14.}
Suponha que $y = y(r)$ seja derivável até a 2\textsuperscript{a} ordem. Verifique que
\[
\dfrac{d}{dr}\!\left[(r^2 + r)\dfrac{dy}{dr}\right] = (2r + 1)\dfrac{dy}{dr} + (r^2 + r)\dfrac{d^2 y}{dr^2}.
\]

\vspace{0.5cm}
\noindent\textbf{Exercício 15.}
Seja $y = x^2$, em que $x = x(t)$ é uma função derivável até a 2\textsuperscript{a} ordem. Verifique que $\dfrac{d^2 y}{dt^2} = 2\!\left(\dfrac{dx}{dt}\right)^{\!2} + 2x\dfrac{d^2 x}{dt^2}$.

\vspace{0.5cm}
\noindent\textbf{Exercício 16.} Suponha que $x = x(t)$ seja derivável até a 2\textsuperscript{a} ordem. Verifique que
\begin{enumerate}
    \renewcommand{\theenumi}{\alph{enumi}}
    \renewcommand{\labelenumi}{\theenumi)}
    \item $\dfrac{d}{dt}\!\left(t^2 \dfrac{dx}{dt}\right) = 2t\dfrac{dx}{dt} + t^2 \dfrac{d^2 x}{dt^2}$
    \item $\dfrac{d}{dt}\!\left(x \dfrac{dx}{dt}\right) = \left(\dfrac{dx}{dt}\right)^{\!2} + x\dfrac{d^2 x}{dt^2}$
\end{enumerate}

\subsection{Regra da Cadeia para Derivação de Função Composta}
Esta subseção não possui exercícios

\subsection{Aplicações da Regra da Cadeia}
\vspace{0.5cm}
\noindent\textbf{Exercício 1.} Determine a derivada.
\begin{enumerate}
    \renewcommand{\theenumi}{\alph{enumi}}
    \renewcommand{\labelenumi}{\theenumi)}
    \item $y = \operatorname{sen} 4x$
    \item $y = \cos 5x$
    \item $f(x) = e^{3x}$
    \item $f(x) = \cos 8x$
    \item $y = \operatorname{sen} t^3$
    \item $g(t) = \ln(2t + 1)$
    \item $x = e^{\operatorname{sen} t}$
    \item $f(x) = \cos e^x$
    \item $y = (\operatorname{sen} x + \cos x)^3$
    \item $y = \sqrt{3x + 1}$
    \item $f(x) = \sqrt[3]{\dfrac{x - 1}{x + 1}}$
    \item $m) \ y = e^{-5x}$
    \item $x = \ln(t^2 + 3t + 9)$
    \item $f(x) = e^{\operatorname{tg} x}$
    \item $y = \operatorname{sen}(\cos x)$
    \item $g(t) = (t^2 + 3)^4$
    \item $f(x) = \cos(x^2 + 3)$
    \item $y = \sqrt{x + e^x}$
    \item $y = \operatorname{tg} 3x$
    \item $y = \sec 3x$
\end{enumerate}

\vspace{0.5cm}
\noindent\textbf{Exercício 2.}
Seja $f\colon \mathbb{R} \to \mathbb{R}$ derivável e seja $g(t) = f(t^2 + 1)$. Supondo $f'(2) = 5$, calcule $g'(1)$.

\vspace{0.5cm}
\noindent\textbf{Exercício 3.}
Seja $f\colon \mathbb{R} \to \mathbb{R}$ derivável e seja $g$ dada por $g(x) = f(e^{2x})$. Supondo $f'(1) = 2$, calcule $g'(0)$.

\vspace{0.5cm}
\noindent\textbf{Exercício 4.} Derive.
\begin{enumerate}
    \renewcommand{\theenumi}{\alph{enumi}}
    \renewcommand{\labelenumi}{\theenumi)}
    \item $y = xe^{3x}$
    \item $y = e^x \cos 2x$
    \item $y = e^{-x}\operatorname{sen} x$
    \item $y = e^{-2t}\operatorname{sen} 3t$
    \item $f(x) = e^{-x^2} + \ln(2x + 1)$
    \item $g(t) = \dfrac{e^t - e^{-t}}{e^t + e^{-t}}$
    \item $y = \dfrac{\cos 5x}{\operatorname{sen} 2x}$
    \item $f(x) = (e^{-x} + e^{x^2})^3$
    \item $y = t^3 e^{-3t}$
    \item $g(x) = e^{x^2}\ln(1 + \sqrt{x})$
    \item $y = (\operatorname{sen} 3x + \cos 2x)^3$
    \item $y = \sqrt{e^x + e^{-x}}$
    \item $y = \ln(x + \sqrt{x^2 + 1})$
    \item $y = \sqrt{x^2 + e^{\sqrt{x}}}$
    \item $y = x\ln(2x + 1)$
    \item $y = [\ln(x^2 + 1)]^3$
    \item $y = \ln(\sec x + \operatorname{tg} x)$
    \item $y = \cos^3 x^3$
    \item $f(x) = \dfrac{\cos x}{\operatorname{sen}^2 x}$
    \item $f(t) = \dfrac{te^{2t}}{\ln(3t + 1)}$
\end{enumerate}

\vspace{0.5cm}
\noindent\textbf{Exercício 5.} Calcule a derivada segunda.
\begin{enumerate}
    \renewcommand{\theenumi}{\alph{enumi}}
    \renewcommand{\labelenumi}{\theenumi)}
    \item $y = \operatorname{sen} 5t$
    \item $y = \cos 4t$
    \item $x = \operatorname{sen}\omega t$, $\omega$ constante
    \item $y = e^{-3x}$
    \item $y = e^{-x^2}$
    \item $y = \dfrac{e^x}{x + 1}$
    \item $y = \ln(x^2 + 1)$
    \item $y = \dfrac{x^2}{x - 1}$
    \item $y = e^{-x} - e^{-2x}$
    \item $y = e^{-x}\cos 2x$
    \item $y = \dfrac{x}{x^2 + 1}$
    \item $y = \dfrac{3x + 1}{x^2 + x}$
    \item $y = \dfrac{\operatorname{sen} 3x}{e^x}$
    \item $y = xe^{-2x}$
    \item $y = \operatorname{sen}(\cos x)$
    \item $f(x) = \dfrac{4x + 5}{x^2 - 1}$
    \item $y = \dfrac{1}{xe^x}$
    \item $y = \dfrac{x^2}{x^2 + x + 1}$
    \item $g(t) = \sqrt{t^2 + 3}$
    \item $y = x^3 \sqrt[3]{x + 2}$
\end{enumerate}

\vspace{0.5cm}
\noindent\textbf{Exercício 6.}
Seja $g\colon \mathbb{R} \to \mathbb{R}$ uma função diferenciável e seja $f$ dada por $f(x) = xg(x^2)$. Verifique que $f'(x) = g(x^2) + 2x^2 g'(x^2)$.

\vspace{0.5cm}
\noindent\textbf{Exercício 7.}
Seja $g\colon \mathbb{R} \to \mathbb{R}$ uma função diferenciável e seja $f$ dada por $f(x) = xg(x^2)$. Calcule $f'(1)$ supondo $g(1) = 4$ e $g'(1) = 2$.

\vspace{0.5cm}
\noindent\textbf{Exercício 8.}
Seja $g\colon \mathbb{R} \to \mathbb{R}$ diferenciável tal que $g(1) = 2$ e $g'(1) = 3$. Calcule $f'(0)$, sendo $f$ dada por $f(x) = e^x g(3x + 1)$.

\vspace{0.5cm}
\noindent\textbf{Exercício 9.}
Seja $f\colon \mathbb{R} \to \mathbb{R}$ derivável até a 2\textsuperscript{a} ordem e seja $g$ dada por $g(x) = f(e^{2x})$. Verifique que $g''(x) = 4e^{2x}[f'(e^{2x}) + e^{2x}f''(e^{2x})]$.

\vspace{0.5cm}
\noindent\textbf{Exercício 10.}
Seja $y = e^{2x}$. Verifique que $\dfrac{d^2 y}{dx^2} - 4y = 0$.

\vspace{0.5cm}
\noindent\textbf{Exercício 11.}
Seja $y = xe^{2x}$. Verifique que $\dfrac{d^2 y}{dx^2} - 4\dfrac{dy}{dx} + 4y = 0$.

\vspace{0.5cm}
\noindent\textbf{Exercício 12.}
Determine $\alpha$ de modo que $y = e^{\alpha x}$ verifique a equação $\dfrac{d^2 y}{dx^2} - 4y = 0$.

\vspace{0.5cm}
\noindent\textbf{Exercício 13.}
Determine $\alpha$ de modo que $y = e^{\alpha x}$ verifique a equação $\dfrac{d^2 y}{dx^2} - 3\dfrac{dy}{dx} + 2y = 0$.

\vspace{0.5cm}
\noindent\textbf{Exercício 14.}
Seja $y = e^{\alpha x}$, em que $\alpha$ é uma raiz da equação $\lambda^2 + a\lambda + b = 0$, com $a$ e $b$ constantes. Verifique que
\[
\dfrac{d^2 y}{dx^2} + a\dfrac{dy}{dx} + by = 0.
\]

\vspace{0.5cm}
\noindent\textbf{Exercício 15.} Seja $g$ uma função derivável. Verifique que
\begin{enumerate}
    \renewcommand{\theenumi}{\alph{enumi}}
    \renewcommand{\labelenumi}{\theenumi)}
    \item $[\operatorname{tg} g(x)]' = \sec^2 g(x) \cdot g'(x)$
    \item $[\sec g(x)]' = \sec g(x)\operatorname{tg} g(x) \cdot g'(x)$
    \item $[\operatorname{cotg} g(x)]' = -\operatorname{cosec}^2 g(x) \cdot g'(x)$
    \item $[\operatorname{cosec} g(x)]' = -\operatorname{cosec} g(x)\operatorname{cotg} g(x) \cdot g'(x)$
\end{enumerate}

\vspace{0.5cm}
\noindent\textbf{Exercício 16.} Derive.
\begin{enumerate}
    \renewcommand{\theenumi}{\alph{enumi}}
    \renewcommand{\labelenumi}{\theenumi)}
    \item $y = \operatorname{tg} 3x$
    \item $y = \sec 4x$
    \item $y = \operatorname{cotg} x^2$
    \item $y = \sec(\operatorname{tg} x)$
    \item $y = \sec x^3$
    \item $y = e^{\operatorname{tg} x^2}$
    \item $y = \operatorname{cosec} 2x$
    \item $y = x^3 \operatorname{tg} 4x$
    \item $y = \ln(\sec 3x + \operatorname{tg} 3x)$
    \item $y = e^{-x}\sec x^2$
    \item $y = (x^2 + \operatorname{cotg} x^2)^3$
    \item $y = x^2 \operatorname{tg} 2x$
\end{enumerate}

\vspace{0.5cm}
\noindent\textbf{Exercício 17.}
Seja $y = \cos \omega t$, $\omega$ constante. Verifique que $\dfrac{d^2 y}{dt^2} + \omega^2 y = 0$.

\vspace{0.5cm}
\noindent\textbf{Exercício 18.}
Seja $y = e^{-t}\cos 2t$. Verifique que $\dfrac{d^2 y}{dt^2} + 2\dfrac{dy}{dt} + 5y = 0$.

\vspace{0.5cm}
\noindent\textbf{Exercício 19.}
Seja $y = \dfrac{x + 1}{x - 1}$. Verifique que $(1 - x)\dfrac{d^2 y}{dx^2} = 2\dfrac{dy}{dx}$.

\vspace{0.5cm}
\noindent\textbf{Exercício 20.}
Seja $y = f(x)$ derivável até a 2\textsuperscript{a} ordem. Verifique que
\[
\dfrac{d}{dx}\!\left(x^2 \dfrac{dy}{dx}\right) = 2x\dfrac{dy}{dx} + x^2 \dfrac{d^2 y}{dx^2}.
\]

\vspace{0.5cm}
\noindent\textbf{Exercício 21.}
Seja $y = \sqrt{x^2 + 1}$. Verifique que $\left(\dfrac{dy}{dx}\right)^{\!2} + y\dfrac{d^2 y}{dx^2} = 1$.

\vspace{0.5cm}
\noindent\textbf{Exercício 22.}
Seja $y = y(x)$ definida no intervalo aberto $I$ e tal que, para todo $x$ em $I$, $\dfrac{dy}{dx} = x^2 + y^2$. Verifique que, para todo $x$ em $I$,
\[
\dfrac{d^2 y}{dx^2} = 2x + 2x^2 y + 2y^3.
\]

\vspace{0.5cm}
\noindent\textbf{Exercício 23.}
Seja $y = f(x)$ uma função derivável num intervalo aberto $I$, com $1 \in I$. Suponha $f(1) = 1$ e que, para todo $x$ em $I$, $f'(x) = x + [f(x)]^3$.
\begin{enumerate}
    \renewcommand{\theenumi}{\alph{enumi}}
    \renewcommand{\labelenumi}{\theenumi)}
    \item Mostre que $f''(x)$ existe para todo $x$ em $I$.
    \item Calcule $f''(1)$.
    \item Determine a equação da reta tangente ao gráfico de $f$ no ponto de abscissa 1.
\end{enumerate}

\vspace{0.5cm}
\noindent\textbf{Exercício 24.}
Seja $y = y(r)$ derivável até a 2\textsuperscript{a} ordem. Verifique que
\[
\dfrac{d}{dr}\!\left(y^2 \dfrac{dy}{dr}\right) = 2y\!\left(\dfrac{dy}{dr}\right)^{\!2} + y^2 \dfrac{d^2 y}{dr^2}.
\]

\vspace{0.5cm}
\noindent\textbf{Exercício 25.}
Seja $y = \dfrac{1}{x^2 + 1}$, em que $x = x(t)$ é uma função definida e derivável em $\mathbb{R}$. Verifique que, para todo $t$ real,
\[
\dfrac{dy}{dt} = -2xy^2 \dfrac{dx}{dt}.
\]

\vspace{0.5cm}
\noindent\textbf{Exercício 26.}
Seja $y = \dfrac{4}{x}$, em que $x = x(t)$ é uma função derivável num intervalo aberto $I$. Suponha que, para todo $t$ em $I$, $x(t) \neq 0$ e $\dfrac{dx}{dt} = \beta$, $\beta$ constante. Verifique que $\dfrac{d^2 y}{dt^2} = \dfrac{8\beta^2}{x^3}$.

\vspace{0.5cm}
\noindent\textbf{Exercício 27.}
Seja $f$ uma função diferenciável e suponha que, para todo $x \in D_f$, $3x^2 + x\operatorname{sen} f(x) = 2$. Mostre que $f'(x) = -\dfrac{6x + \operatorname{sen} f(x)}{x\cos f(x)}$, para todo $x \in D_f$, com $x\cos f(x) \neq 0$.

\vspace{0.5cm}
\noindent\textbf{Exercício 28.}
A função diferenciável $y = f(x)$ é tal que, para todo $x \in D_f$, o ponto $(x, f(x))$ é solução da equação $xy^3 + 2xy^2 + x = 4$. Sabe-se que $f(1) = 1$. Calcule $f'(1)$.

\vspace{0.5cm}
\noindent\textbf{Exercício 29.} Seja $f\colon {]-r, r[} \to \mathbb{R}$ uma função derivável. Prove
\begin{enumerate}
    \renewcommand{\theenumi}{\alph{enumi}}
    \renewcommand{\labelenumi}{\theenumi)}
    \item Se $f$ for uma função ímpar, então $f'$ será par.
    \item Se $f$ for função par, então $f'$ será ímpar.
\end{enumerate}

\vspace{0.5cm}
\noindent\textbf{Exercício 30.}
Seja $g\colon \mathbb{R} \to \mathbb{R}$ uma função diferenciável tal que $g(2) = 2$ e $g'(2) = 2$. Calcule $H'(2)$, sendo $H$ dada por $H(x) = g(g(g(x)))$.

\subsection{Derivada de $f(x)^{g(x)}$}
\vspace{0.5cm}
\noindent\textbf{Exercício 1.} Calcule a derivada.
\begin{enumerate}
    \renewcommand{\theenumi}{\alph{enumi}}
    \renewcommand{\labelenumi}{\theenumi)}
    \item $f(x) = 5^x + \log_3 x$
    \item $y = 2^{x^2} + 3^{2x}$
    \item $g(x) = 3^{2x + 1} + \log_2 (x^2 + 1)$
    \item $y = (2x + 1)^x$
    \item $f(x) = x^{\operatorname{sen} 3x}$
    \item $g(x) = (3 + \cos x)^x$
    \item $y = x^x \operatorname{sen} x$
    \item $y = x^{x^2 + 1}$
    \item $y = (1 + i)^{-t}$, $i$ constante
    \item $y = 10^x - 10^{-x}$
    \item $y = (2 + \operatorname{sen} x)^{\cos 3x}$
    \item $y = \ln(1 + x^x)$
    \item $y = \left(1 + \dfrac{1}{x}\right)^{\!x}$
    \item $y = x^{x^x}$
    \item $y = x^\pi + \pi^x$
    \item $y = (1 + x)^{e^{-x}}$
\end{enumerate}

\vspace{0.5cm}
\noindent\textbf{Exercício 2.}
Sejam $f$ e $g$ deriváveis em $A$, com $f(x) > 0$ em $A$. Verifique que, para todo $x$ em $A$,
\[
[f(x)^{g(x)}]' = f(x)^{g(x)} g'(x) \ln f(x) + g(x) f(x)^{g(x) - 1} f'(x)
\]

\vspace{0.5cm}
\noindent\textbf{Exercício 3.}
Utilizando o resultado obtido no Exercício 2, calcule a derivada.
\begin{enumerate}
    \renewcommand{\theenumi}{\alph{enumi}}
    \renewcommand{\labelenumi}{\theenumi)}
    \item $y = (x + 2)^x$
    \item $y = (1 + e^x)^{x^2}$
    \item $y = (4 + \operatorname{sen} 3x)^x$
    \item $y = (x + 3)^{x^2}$
    \item $y = (3 + \pi)^{x^2}$
    \item $y = (x^2 + 1)^\pi$
\end{enumerate}

\subsection{Derivação da Função Dada Implicitamente}
\vspace{0.5cm}
\noindent\textbf{Exercício 1.}
Suponha que $y = f(x)$ seja uma função derivável e dada implicitamente pela equação
\[
xy^2 + y + x = 1.
\]
Mostre que $f'(x) = \dfrac{-1 - [f(x)]^2}{2xf(x) + 1}$ em todo $x \in D_f$ com $2xf(x) + 1 \neq 0$.

\vspace{0.5cm}
\noindent\textbf{Exercício 2.}
Determine uma função $y = f(x)$ que seja dada implicitamente pela equação $xy^2 + y + x = 1$.

\vspace{0.5cm}
\noindent\textbf{Exercício 3.}
A função $y = f(x)$ é dada implicitamente pela equação $xy + 3 = 2x$. Mostre que $x\dfrac{dy}{dx} = 2 - y$. Calcule $\left.\dfrac{dy}{dx}\right|_{x = 2}$.

\vspace{0.5cm}
\noindent\textbf{Exercício 4.}
Expresse $\dfrac{dy}{dx}$ em termos de $x$ e de $y$, em que $y = f(x)$ é uma função diferenciável dada implicitamente pela equação
\begin{enumerate}
    \renewcommand{\theenumi}{\alph{enumi}}
    \renewcommand{\labelenumi}{\theenumi)}
    \item $x^2 - y^2 = 4$
    \item $y^3 + x^2 y = x + 4$
    \item $xy^2 + 2y = 3$
    \item $y^5 + y = x$
    \item $x^2 + 4y^2 = 3$
    \item $xy + y^3 = x$
    \item $x^2 + y^2 + 2y = 0$
    \item $xy^3 + xy = 2$
    \item $xe^y + xy = 3$
    \item $y + \ln(x^2 + y^2) = 4$
    \item $5y + \cos y = xy$
    \item $2y + \operatorname{sen} y = x$
\end{enumerate}

\vspace{0.5cm}
\noindent\textbf{Exercício 5.}
A função $y = f(x)$, $y > 0$, é dada implicitamente por $x^2 + 4y^2 = 2$. Determine a equação da reta tangente ao gráfico de $f$, no ponto de abscissa 1.

\vspace{0.5cm}
\noindent\textbf{Exercício 6.}
Determine a equação da reta tangente à elipse $\dfrac{x^2}{a^2} + \dfrac{y^2}{b^2} = 1$, no ponto $(x_0, y_0)$, $y_0 \neq 0$.

\vspace{0.5cm}
\noindent\textbf{Exercício 7.}
Verifique que $y_0 x + x_0 y = 2$ é a equação da reta tangente à curva $xy = 1$ no ponto $(x_0, y_0)$, $x_0 > 0$. Conclua que $(x_0, y_0)$ é o ponto médio do segmento $AB$, em que $A$ e $B$ são as interseções da reta tangente, em $(x_0, y_0)$, com os eixos coordenados.

\vspace{0.5cm}
\noindent\textbf{Exercício 8.}
Suponha que $y = f(x)$ seja uma função derivável dada implicitamente pela equação $y^3 + 2xy^2 + x = 4$. Suponha, ainda, que $1 \in D_f$.
\begin{enumerate}
    \renewcommand{\theenumi}{\alph{enumi}}
    \renewcommand{\labelenumi}{\theenumi)}
    \item Calcule $f(1)$.
    \item Determine a equação da reta tangente ao gráfico de $f$ no ponto de abscissa 1.
\end{enumerate}

\vspace{0.5cm}
\noindent\textbf{Exercício 9.}
A reta tangente à curva $x^{\frac{2}{3}} + y^{\frac{2}{3}} = 1$, no ponto $(x_0, y_0)$, $x_0 > 0$ e $y_0 > 0$, intercepta os eixos $x$ e $y$ nos pontos $A$ e $B$, respectivamente. Mostre que a distância de $A$ e $B$ não depende de $(x_0, y_0)$.

\vspace{0.5cm}
\noindent\textbf{Exercício 10.}
A reta tangente à curva $xy - x^2 = 1$ no ponto $(x_0, y_0)$, $x_0 > 0$, intercepta o eixo $y$ no ponto $B$. Mostre que a área do triângulo de vértices $(0,0)$, $(x_0, y_0)$ e $B$ não depende de $(x_0, y_0)$.

\vspace{0.5cm}
\noindent\textbf{Exercício 11.}
A função $y = f(x)$ é dada implicitamente pela equação $3y^2 + 2xy - x^2 = 3$. Sabe-se que, para todo $x \in D_f$, $f(x) > 0$ e que $f$ admite uma reta tangente $T$ paralela à reta $5y - x = 2$. Determine $T$.

\subsection{Interpretação de $\frac{dy}{dx}$ como um Quociente Diferencial}
\vspace{0.5cm}
\noindent\textbf{Exercício 1.} Calcule a diferencial.
\begin{enumerate}
    \renewcommand{\theenumi}{\alph{enumi}}
    \renewcommand{\labelenumi}{\theenumi)}
    \item $y = x^3$
    \item $y = x^2 - 2x$
    \item $y = \dfrac{x}{x + 1}$
    \item $y = \sqrt[3]{x}$
\end{enumerate}

\vspace{0.5cm}
\noindent\textbf{Exercício 2.} Seja $A = l^2$, $l > 0$.
\begin{enumerate}
    \renewcommand{\theenumi}{\alph{enumi}}
    \renewcommand{\labelenumi}{\theenumi)}
    \item Calcule a diferencial.
    \item Interprete geometricamente o erro que se comete na aproximação de $\Delta A$ por $dA$. (Olhe para $A = l^2$ como a fórmula para o cálculo da área do quadrado de lado $l$.)
\end{enumerate}

\vspace{0.5cm}
\noindent\textbf{Exercício 3.} Seja $V = \dfrac{4}{3}\pi r^3$, $r > 0$.
\begin{enumerate}
    \renewcommand{\theenumi}{\alph{enumi}}
    \renewcommand{\labelenumi}{\theenumi)}
    \item Calcule a diferencial.
    \item Interprete geometricamente $dV$. (Lembre-se de que $V$ é o volume da esfera de raio $r$ e que $4\pi r^2$ é a área da superfície esférica de raio $r$.)
\end{enumerate}

\vspace{0.5cm}
\noindent\textbf{Exercício 4.} Seja $y = x^2 + 3x$.
\begin{enumerate}
    \renewcommand{\theenumi}{\alph{enumi}}
    \renewcommand{\labelenumi}{\theenumi)}
    \item Calcule a diferencial.
    \item Calcule o erro que se comete na aproximação de $\Delta y$ por $dy$. Interprete graficamente.
\end{enumerate}

\subsection{Velocidade e Aceleração. Taxa de Variação}
\vspace{0.5cm}
\noindent\textbf{Exercício 1.} Uma partícula desloca-se sobre o eixo $x$ com função de posição $x = 3 + 2t - t^2$, $t \ge 0$.
\begin{enumerate}
    \renewcommand{\theenumi}{\alph{enumi}}
    \renewcommand{\labelenumi}{\theenumi)}
    \item Qual a velocidade no instante $t$?
    \item Qual a aceleração no instante $t$?
    \item Estude a variação do sinal de $v(t)$.
    \item Esboce o gráfico da função de posição.
\end{enumerate}

\vspace{0.5cm}
\noindent\textbf{Exercício 2.} Uma partícula desloca-se sobre o eixo $x$ com função de posição $x = \dfrac{1}{2}t + 1$, $t \ge 0$.
\begin{enumerate}
    \renewcommand{\theenumi}{\alph{enumi}}
    \renewcommand{\labelenumi}{\theenumi)}
    \item Determine a velocidade no instante $t$.
    \item Qual a aceleração no instante $t$?
    \item Esboce o gráfico da função de posição.
\end{enumerate}

\vspace{0.5cm}
\noindent\textbf{Exercício 3.} A posição de uma partícula que se desloca ao longo do eixo $x$ depende do tempo de acordo com a equação $x = -t^3 + 3t^2$, $t \ge 0$.
\begin{enumerate}
    \renewcommand{\theenumi}{\alph{enumi}}
    \renewcommand{\labelenumi}{\theenumi)}
    \item Estude o sinal de $v(t)$.
    \item Estude o sinal de $a(t)$.
    \item Calcule $\lim\limits_{t \to +\infty} (-t^3 + 3t^2)$.
    \item Esboce o gráfico da função $x = -t^3 + 3t^2$, $t \ge 0$.
\end{enumerate}

\vspace{0.5cm}
\noindent\textbf{Exercício 4.}
Seja $x = f(t)$, $t \ge 0$, tal que $f(0) = 1$, $\dfrac{dx}{dt} > 0$ e $\dfrac{d^2 x}{dt^2} > 0$ para $t \ge 0$. Como você acha que deve ser o gráfico de $f$? Por quê?

\vspace{0.5cm}
\noindent\textbf{Exercício 5.}
A função $x = f(t)$, $t \in I$, é derivável até a 2\textsuperscript{a} ordem no intervalo aberto $I$ e seu gráfico tem o seguinte aspecto

\begin{center}
\begin{tikzpicture}[scale=0.9, >=latex]
    % Eixos
    \draw[->, thick] (-1,0) -- (4.5,0) node[below] {$t$};
    \draw[->, thick] (0,-0.5) -- (0,3.5) node[left] {$x$};
    \node[below left] at (0,0) {$0$};

    % Curva f côncava para baixo
    \draw[ultra thick] (0.5, 0.8) to[out=60, in=175] (3.5, 2.7) node[above right] {$f$};

    % Linhas tracejadas
    \draw[dashed] (0.5, 0) -- (0.5, 0.8);
    \draw[dashed] (3.5, 0) -- (3.5, 2.7);

    % Intervalo I no eixo t
    \draw[ultra thick] (0.5, 0) -- (3.5, 0);
    \node[below=4pt] at (2,0) {$I$};
\end{tikzpicture}
\end{center}

\noindent O que é mais razoável esperar que ocorra: $f''(t) \le 0$ ou $f''(t) > 0$ em $I$? Por quê?

\vspace{0.5cm}
\noindent\textbf{Exercício 6.}
Seja $x = f(t)$, $t \ge 0$, tal que $f(0) = 1$ e $f(1) = 2$. Suponha, ainda, que $\dfrac{dx}{dt} > 0$ para $t \ge 0$, $\dfrac{d^2x}{dt^2} < 0$ para $0 < t < 1$ e $\dfrac{d^2x}{dt^2} > 0$ para $t > 1$. Como você acha que deve ser o gráfico de $f$? Por quê?

\vspace{0.5cm}
\noindent\textbf{Exercício 7.} Seja $f(t) = t^3 + 3t^2$.
\begin{enumerate}
    \renewcommand{\theenumi}{\alph{enumi}}
    \renewcommand{\labelenumi}{\theenumi)}
    \item Estude o sinal de $f'(t)$.
    \item Estude o sinal de $f''(t)$.
    \item Calcule $\lim\limits_{t \to +\infty} (t^3 + 3t^2)$ e $\lim\limits_{t \to -\infty} (t^3 + 3t^2)$.
    \item Utilizando as informações acima, esboce o gráfico de $f$.
\end{enumerate}

\vspace{0.5cm}
\noindent\textbf{Exercício 8.} Seja $f(t) = \dfrac{t}{t^2 + 4}$.
\begin{enumerate}
    \renewcommand{\theenumi}{\alph{enumi}}
    \renewcommand{\labelenumi}{\theenumi)}
    \item Estude o sinal de $f'(t)$.
    \item Estude o sinal de $f''(t)$.
    \item Calcule $\lim\limits_{t \to +\infty} \dfrac{t}{t^2 + 4}$ e $\lim\limits_{t \to -\infty} \dfrac{t}{t^2 + 4}$.
    \item Utilizando as informações acima, esboce o gráfico de $f$.
\end{enumerate}

\vspace{0.5cm}
\noindent\textbf{Exercício 9.} A posição de uma partícula que se desloca ao longo do eixo $x$ varia com o tempo segundo a equação $x = \dfrac{v_0}{k}(1 - e^{-kt})$, $t \ge 0$, em que $v_0$ e $k$ são constantes estritamente positivas.
\begin{enumerate}
    \renewcommand{\theenumi}{\alph{enumi}}
    \renewcommand{\labelenumi}{\theenumi)}
    \item Qual a velocidade no instante $t$?
    \item Com argumentos físicos, justifique a afirmação: ``a função é estritamente crescente.''
    \item Qual a aceleração no instante $t$?
    \item Com argumentos físicos, justifique a afirmação: ``o gráfico da função tem a concavidade voltada para baixo.''
    \item Calcule $\lim\limits_{t \to +\infty} \dfrac{v_0}{k}(1 - e^{-kt})$.
    \item Esboce o gráfico da função.
\end{enumerate}

\vspace{0.5cm}
\noindent\textbf{Exercício 10.} A equação do movimento de uma partícula que se desloca ao longo do eixo $x$ é $x = e^{-t}\operatorname{sen} t$, $t \ge 0$.
\begin{enumerate}
    \renewcommand{\theenumi}{\alph{enumi}}
    \renewcommand{\labelenumi}{\theenumi)}
    \item Determine a velocidade e a aceleração no instante $t$.
    \item Calcule $\lim\limits_{t \to +\infty} e^{-t}\operatorname{sen} t$.
    \item Esboce o gráfico da função.
    \item Interprete tal movimento.
\end{enumerate}

\vspace{0.5cm}
\noindent\textbf{Exercício 11.}
Um ponto $P$ move-se sobre a parábola $y = 3x^2 - 2x$. Suponha que as coordenadas $x(t)$ e $y(t)$ de $P$ são deriváveis e que $\dfrac{dx}{dt} \neq 0$. Pergunta-se: em que ponto da parábola a velocidade da ordenada $y$ de $P$ é o triplo da velocidade da abscissa $x$ de $P$?

\vspace{0.5cm}
\noindent\textbf{Exercício 12.}
Um ponto $P$ move-se ao longo do gráfico de $y = \dfrac{1}{x^2 + 1}$ de tal modo que a sua abscissa $x$ varia a uma velocidade constante de $5\text{ (m/s)}$. Qual a velocidade de $y$ no instante em que $x = 10\text{ m}$?

\vspace{0.5cm}
\noindent\textbf{Exercício 13.}
Um ponto desloca-se sobre a hipérbole $xy = 4$ de tal modo que a velocidade de $y$ é $\dfrac{dy}{dt} = \beta$, $\beta$ constante. Mostre que a aceleração da abscissa $x$ é $\dfrac{d^2x}{dt^2} = \dfrac{\beta^2}{8}x^3$.

\vspace{0.5cm}
\noindent\textbf{Exercício 14.} Seja um ponto move-se ao longo da elipse $x^2 + 4y^2 = 1$. A abscissa $x$ está variando a uma velocidade $\dfrac{dx}{dt} = \operatorname{sen} 4t$. Mostre que
\begin{enumerate}
    \renewcommand{\theenumi}{\alph{enumi}}
    \renewcommand{\labelenumi}{\theenumi)}
    \item $\dfrac{dy}{dt} = -\dfrac{x\operatorname{sen} 4t}{4y}$
    \item $\dfrac{d^2y}{dt^2} = -\dfrac{\operatorname{sen}^2 4t + 16xy^2 \cos 4t}{16y^3}$
\end{enumerate}

\vspace{0.5cm}
\noindent\textbf{Exercício 15.}
Um ponto move-se sobre a semicircunferência $x^2 + y^2 = 5$, $y \ge 0$. Suponha $\dfrac{dx}{dt} > 0$. Determine o ponto da curva em que a velocidade de $y$ seja o dobro da de $x$.

\vspace{0.5cm}
\noindent\textbf{Exercício 16.}
Uma escada de $8\text{ m}$ está encostada em uma parede. Se a extremidade inferior da escada for afastada do pé da parede a uma velocidade constante de $2\text{ (m/s)}$, com que velocidade a extremidade superior estará descendo no instante em que a inferior estiver a $3\text{ m}$ da parede?

\vspace{0.5cm}
\noindent\textbf{Exercício 17.}
Suponha que os comprimentos dos segmentos $AB$ e $OB$ sejam, respectivamente, $5\text{ cm}$ e $3\text{ cm}$. Suponha, ainda, que $\theta$ esteja variando a uma taxa constante de $\dfrac{1}{2}\text{ rad/s}$. Determine a velocidade de $A$, quando $\theta = \dfrac{\pi}{2}\text{ rad}$.
\begin{center}
\begin{tikzpicture}[scale=0.9, >=latex]
    % Eixos coordenados
    \draw[->, thick] (-0.8, 0) -- (7.5, 0) node[below] {$x$};
    \draw[->, thick] (0, -0.8) -- (0, 4.2) node[left] {$y$};
    \node[below left] at (0, 0) {$O$};
    
    % Vértices
    \coordinate (O) at (0, 0);
    \coordinate (B) at (1.5, 2.6);
    \coordinate (A) at (5.8, 0);
    
    % Segmentos do mecanismo/triângulo
    \draw[line width=1.5pt, blue!60!black] (O) -- (B) node[midway, above left, text=black] {$3$};
    \draw[line width=1.5pt, blue!60!black] (B) -- (A) node[midway, above right, text=black] {$5$};
    
    % Rótulos dos pontos
    \node[above] at (B) {$B$};
    \node[above right=1pt] at (A) {$A$};
    \node[below=3pt] at (A) {$(x, 0)$};
    
    % Ângulo theta
    \draw[->] (1.1, 0) arc (0:60:1.1);
    \node at (1.4, 0.45) {$\theta$};
\end{tikzpicture}
\end{center}

\vspace{0.5cm}
\noindent\textbf{Exercício 18.}
Enche-se um reservatório, cuja forma é a de um cone circular reto, de água a uma taxa de $0{,}1\text{ m}^3\text{/s}$. O vértice está a $15\text{ m}$ do topo e o raio do topo é de $10\text{ m}$. Com que velocidade o nível $h$ da água está subindo no instante em que $h = 5\text{ m}$.

\begin{center}
\begin{tikzpicture}[scale=0.35]
    % Elipse superior do cone
    \draw[thick] (0,15) ellipse (10 and 2.5);
    
    % Paredes laterais do cone
    \draw[thick] (-10,15) -- (0,0) -- (10,15);
    
    % Nível da água (h = 5, raio = 10/3 ~ 3.33)
    \draw[thick, dashed] (0,5) ellipse (3.33 and 0.8);
    \fill[gray!20, opacity=0.6] (-3.33,5) -- (0,0) -- (3.33,5) arc (0:-180:3.33 and 0.8) -- cycle;
    \draw[thick] (-3.33,5) arc (-180:0:3.33 and 0.8);
    
    % Eixo central e cotas verticais
    \draw[dashed] (0,15) -- (0,0);
    \draw[<->] (-12,0) -- (-12,15) node[midway, fill=white, inner sep=2pt] {$15\text{ m}$};
    \draw[<->] (5,0) -- (5,5) node[midway, right] {$h$};
    \draw[<->] (0,15) -- (10,15) node[midway, above] {$10\text{ m}$};
\end{tikzpicture}
\end{center}

\vspace{0.5cm}
\noindent\textbf{Exercício 19.}
O ponto $P = (x, y)$ está fixo à roda de raio $1\text{ m}$, que rola, sem escorregamento, sobre o eixo $x$. O ângulo $\theta$ está variando a uma taxa constante de $1\text{ rad/s}$. Expresse as velocidades da abscissa e da ordenada de $P$ em função de $\theta$.
\begin{center}
\begin{tikzpicture}[scale=1.1, >=latex]
    % Eixos coordenados
    \draw[->, thick] (-0.8, 0) -- (5.5, 0) node[below] {$x$};
    \draw[->, thick] (0, -0.6) -- (0, 3.8) node[left] {$y$};
    
    % Centro e círculo apoiado no eixo x
    \coordinate (C) at (2.4, 1.6);
    \draw[line width=2pt, blue!50!black] (C) circle (1.6);
    \node[above=1pt] at (C) {$0$};
    
    % Ponto de tangência na base
    \coordinate (T) at (2.4, 0);
    
    % Linha vertical tracejada do centro até a base
    \draw[dashed, thick] (C) -- (T);
    
    % Ponto P na borda do círculo (coordenada polar pura, sem depender da lib calc)
    \path (C) ++(210:1.6) coordinate (P);
    
    % Raio tracejado até o ponto P
    \draw[dashed, thick] (C) -- (P);
    \node[left=2pt] at (P) {$P$};
    
    % Arco e rótulo do ângulo teta medido em relação à vertical descendente
    \draw[->] (2.4, 0.9) arc (-90:-150:0.7);
    \node at (2.05, 0.65) {$\theta$};
\end{tikzpicture}
\end{center}

\vspace{0.5cm}
\noindent\textbf{Exercício 20.}
Um ponto $P$ move-se sobre a parábola $y^2 = x$, $x > 0$ e $y > 0$. A abscissa $x$ está variando com uma aceleração que, em cada instante, é o dobro do quadrado da velocidade da ordenada $y$. Mostre que a ordenada está variando com aceleração nula.

\vspace{0.5cm}
\noindent\textbf{Exercício 21.}
Dois pontos $P$ e $Q$ deslocam-se, respectivamente, nos eixos $x$ e $y$ de modo que a soma das distâncias de $P$ a $R$ e de $R$ a $Q$ mantém-se constante e igual a $e$ durante o movimento, em que $R = (0, h)$ é um ponto fixo. (Veja a figura a seguir.)

\begin{center}
\begin{tikzpicture}[scale=0.9, >=latex]
    % Eixos
    \draw[->, thick] (-1,0) -- (5.5,0) node[below] {$x$};
    \draw[->, thick] (0,-0.5) -- (0,3.5) node[left] {$y$};
    \node[below left] at (0,0) {$0$};
    
    % Pontos
    \coordinate (R) at (0,2.3);
    \coordinate (Q) at (0,0.8);
    \coordinate (P) at (4,0);
    
    % Segmentos da corda/haste
    \draw[ultra thick] (Q) -- (R) -- (P);
    
    % Pontos marcados
    \fill (R) circle (2pt) node[above right] {$R$};
    \fill (Q) circle (2pt) node[left] {$Q$} node[right=4pt] {$y$};
    \fill (P) circle (2pt) node[above right] {$P$} node[below=4pt] {$x$};
    
    % Linha tracejada da cota h
    \draw[dashed] (-1.2, 2.3) -- (0, 2.3);
    \draw[dashed] (-1.2, 0) -- (0, 0);
    \draw[<->] (-1, 0) -- (-1, 2.3) node[midway, fill=white, inner sep=1pt] {$h$};
\end{tikzpicture}
\end{center}

\noindent Relacione a velocidade $\dfrac{dy}{dt}$ de $Q$ com a velocidade $\dfrac{dx}{dt}$ de $P$.

\subsection{Problemas Reta Tangente e Normal ao Gráfico de uma Função}
\subsection{Exercícios do Capítulo}

\end{document}
